% Capítulo "Família e Cuidados" do relatório -- EDITÁVEL À MÃO (specs/2026-11-05_versao_final_site R5).
% Exportado uma vez por gera_latex.py --exportar-capitulos (2026-11-05); desde então o TEXTO é escrito aqui, não
% vem mais de relatorio/textos_curados.json (que é só do site). Os DADOS continuam gerados: as figuras
% (\figura{...}{../../visualizacoes/a4/...} e mapas/a4/) saem do analise.py; as tabelas no corpo
% (\input{gerado/tabelas_corpo/...}) e os rótulos do apêndice (\ref{tab:...}) saem do gera_latex.py, que confere
% a cada execução se tudo o que este arquivo cita existe.

\chapter{Família e Cuidados}\label{cap:eixo-3}

\begin{achados}\begin{itemize}

\item Há tendência de redução no número de matrículas escolares

\item Mesmo assim, o atendimento segue abaixo das metas do Plano Nacional de Educação: 85,2\% na pré-escola (4 e 5 anos, meta de 100\%) e 43,7\% na creche (0 a 3 anos, meta de 50\%) em 2025

\item Predominam arranjos familiares com uma única mulher adulta (dado que precisa ser tomado com cautela frente a natureza administrativa dos dados)

\end{itemize}\end{achados}

A natureza dos dados administrativos, especialmente os advindos do Cadastro Único, exige cautela na análise. Dito isso, a predominância de arranjos familiares com apenas uma mulher adulta exige atenção. Já os dados educacionais apontam taxas de atendimento abaixo das metas do Plano Nacional de Educação, com a maior distância na pré-escola: 85,2\% das crianças de 4 e 5 anos em 2025, para uma meta de 100\%; na creche, 43,7\% das crianças de 0 a 3 anos, para uma meta de 50\%. A oferta de matrículas da rede pública precisará também se atentar a mudança na demanda geográfica.

\section{Crianças até 72 meses frequentando escola/\allowbreak{}creche (geral)}

\figura{grafico}{../../visualizacoes/a4/sidra_frequencia_escola_0_5_total_2022.pdf}{Crianças de 0 a 5 anos que frequentam escola/creche, por idade (Censo 2022)}{Elaboração IPP com dados de \citeonline{ibge_censo2022}.}{graf:sidra-frequencia-escola-0-5-total-2022}

Em 2022, 233.509 crianças de 0 a 5 anos frequentavam escola ou creche no município. As crianças de 4 e 5 anos somavam 127.369, enquanto as de 0 a 3 anos totalizavam 106.140. A diferença entre os grupos etários permite analisar como a frequência escolar se distribui ao longo da primeira infância e estabelecer comparações com os dados de matrículas. O indicador também pode ser relacionado à população de cada faixa etária, permitindo observar a participação escolar em relação ao contingente infantil.

\vertabelas{Tabela~\ref{tab:sidra-frequencia-escola-0-5-total-2022}; ver Apêndice~\ref{ap:eixo-3}}

\section{Crianças até 72 meses frequentando escola/\allowbreak{}creche, por raça/\allowbreak{}cor}

\figura{grafico}{../../visualizacoes/a4/sidra_taxa_frequencia_0_6_raca_2022.pdf}{Taxa de frequência escolar bruta (0 a 5 anos), por idade e raça/cor (Censo 2022)}{Elaboração IPP com dados de \citeonline{ibge_censo2022}. Nota: Amarela e indígena (grupos pequenos) só no total de 0 a 5 anos: 62,8\%.}{graf:sidra-taxa-frequencia-0-6-raca-2022}

A taxa de frequência escolar bruta aumenta conforme a idade, passando de 8,04\% entre crianças de 0 ano para 90,57\% aos 5 anos. No conjunto de 0 a 5 anos, a taxa foi de 59,30\%, com diferenças entre os grupos de raça/cor: 62,26\% entre crianças pretas, 59,10\% entre pardas e 58,70\% entre brancas. A comparação por idade permite analisar como a frequência escolar se modifica ao longo da primeira infância e como esse comportamento varia entre os grupos de raça/cor. Os dados podem ser relacionados ao número absoluto de crianças frequentando escola/creche para complementar a análise.

\vertabelas{Tabela~\ref{tab:sidra-taxa-frequencia-0-6-raca-2022}; ver Apêndice~\ref{ap:eixo-3}}

\section{Crianças até 72 meses frequentando escola/\allowbreak{}creche, por sexo}

\figura{grafico}{../../visualizacoes/a4/sidra_taxa_frequencia_0_6_sexo_2022.pdf}{Taxa de frequência escolar bruta (0 a 5 anos), por idade e sexo (Censo 2022)}{Elaboração IPP com dados de \citeonline{ibge_censo2022}.}{graf:sidra-taxa-frequencia-0-6-sexo-2022}

A taxa de frequência escolar bruta aumenta conforme a idade, passando de 8,04\% aos 0 anos para 90,57\% aos 5 anos. No conjunto de 0 a 5 anos, a taxa foi de 59,81\% entre os meninos e 58,78\% entre as meninas. A diferença entre os sexos varia ao longo das idades: aos 4 anos, a taxa foi de 82,51\% entre meninos e 83,35\% entre meninas, enquanto aos 5 anos foi de 91,56\% entre meninos e 89,49\% entre meninas. A comparação por idade e sexo permite analisar como a frequência escolar se modifica ao longo da primeira infância.

\vertabelas{Tabela~\ref{tab:sidra-taxa-frequencia-0-6-sexo-2022}; ver Apêndice~\ref{ap:eixo-3}}

\section{Taxa bruta de frequência escolar da população até 72 meses}

\figura{grafico}{../../visualizacoes/a4/sidra_taxa_frequencia_0_5_total_2022.pdf}{Taxa de frequência escolar bruta por idade, 0 a 5 anos (Censo 2022)}{Elaboração IPP com dados de \citeonline{ibge_censo2022}.}{graf:sidra-taxa-frequencia-0-5-total-2022}

A frequência escolar na primeira infância apresenta uma trajetória de crescimento acelerado à medida que a idade da criança vai aumentando. Esse movimento pode ser explicado pela necessidade de retorno dos pais, em especial das mães, ao mercado de trabalho e garantia do direito constitucional ao desenvolvimento para as crianças. A partir dos 4 anos, quando há a obrigatoriedade legal da pré-escola a taxa sobe para cerca de 83\%, atingindo 90\% aos 5 anos. A despeito do alto percentual, é um ponto de atenção ter uma déficit de 17\% e 10\% de crianças em idade escolar obrigatória que não a estejam frequentando.

\vertabelas{Tabela~\ref{tab:sidra-taxa-frequencia-0-5-total-2022}; ver Apêndice~\ref{ap:eixo-3}}

\section{Matrículas na educação básica de crianças até 72 meses}

\figura{grafico}{../../visualizacoes/a4/matriculas_0_a_5_por_ano.pdf}{Matrículas de crianças de 0 a 5 anos por ano}{Elaboração IPP com dados de \citeonline{inep_censo_escolar}.}{graf:matriculas-0-a-5-por-ano}

O número de matrículas de crianças de 0 a 5 anos passou de 170.650 em 2007 para 258.868 em 2019, crescimento de aproximadamente 52\%. A partir de 2020, observa-se uma redução, com 230.260 matrículas em 2021, seguida de recuperação parcial em 2022. Em 2025, foram registradas 230.284 matrículas, cerca de 11\% abaixo do pico de 2019. A série permite relacionar a evolução das matrículas às mudanças na população de 0 a 5 anos e comparar períodos de expansão e redução da participação escolar.

\figura{grafico}{../../visualizacoes/a4/matriculas_0_a_5_creche_pre_por_ano.pdf}{Matrículas de crianças de 0 a 5 anos, por faixa de idade}{Elaboração IPP com dados de \citeonline{inep_censo_escolar}.}{graf:matriculas-0-a-5-creche-pre-por-ano}

As matrículas de crianças de 0 a 3 anos passaram de 60.997 em 2007 para 115.723 em 2019, enquanto as de 4 a 5 anos aumentaram de 109.653 para 143.145 no mesmo período. Após 2019, os dois grupos apresentam redução, chegando em 2025 a 110.097 matrículas de 0 a 3 anos e 120.187 de 4 a 5 anos. A série permite comparar a evolução das matrículas entre os dois recortes etários e observar diferenças na trajetória de cada grupo ao longo do período analisado.

\figura{grafico}{../../visualizacoes/a4/matriculas_0_a_5_rede_por_ano.pdf}{Matrículas de crianças de 0 a 5 anos, por rede}{Elaboração IPP com dados de \citeonline{inep_censo_escolar}.}{graf:matriculas-0-a-5-rede-por-ano}

As matrículas na rede pública passaram de 107.599 em 2007 para 144.314 em 2019, enquanto na rede privada aumentaram de 63.051 para 114.554 no mesmo período. Após 2019, a rede pública apresentou redução contínua, chegando a 122.236 matrículas em 2025. Na rede privada, houve queda até 2021, seguida de recuperação parcial, com 108.048 matrículas em 2025. A evolução dos dois segmentos permite acompanhar mudanças na participação de cada rede no atendimento às crianças de 0 a 5 anos ao longo da série.

\vertabelas{Tabela~\ref{tab:matriculas-0-a-5-por-ano}; ver Apêndice~\ref{ap:eixo-3}}

\section{Taxa bruta de atendimento escolar até 72 meses}

\figura{grafico}{../../visualizacoes/a4/taxa_atendimento_0_a_5_por_ano.pdf}{Taxa bruta de atendimento escolar de 0 a 5 anos (\%)}{Elaboração IPP com dados de \citeonline{ms_ripsa_populacao}; \citeonline{inep_censo_escolar}.}{graf:taxa-atendimento-0-a-5-por-ano}

A taxa bruta de atendimento escolar de crianças de 0 a 5 anos passou de 34,4\% em 2007 para 58,6\% em 2025. No mesmo período, a taxa entre crianças de 0 a 3 anos aumentou de 18,9\% para 43,7\%, enquanto entre 4 e 5 anos passou de 63,4\% para 85,2\%. A diferença entre os grupos etários permanece ao longo da série, permitindo comparar os níveis de atendimento e sua evolução segundo a idade. O indicador pode ser relacionado às matrículas e à população de cada faixa etária para analisar a cobertura do atendimento escolar.

\vertabelas{Tabela~\ref{tab:matriculas-0-a-5-por-ano}; ver Apêndice~\ref{ap:eixo-3}}

\section{Cobertura vacinal de rotina em crianças até 2 anos}

\figura{grafico}{../../visualizacoes/a4/cobertura_vacinal_epi_ano.pdf}{Cobertura vacinal por imunobiológico (2016-2026)}{Elaboração IPP com dados de \citeonline{sms_rio_epi_vacinal}.}{graf:cobertura-vacinal-epi-ano}

A evolução da cobertura vacinal na cidade do Rio de Janeiro traz uma trajetória com uma elevada cobertura para grande parte dos imunizantes até 2018, momento que se inicia uma redução entre os anos de 2019 e 2022. A partir de 2023 observa-se recuperação das coberturas, especialmente para pneumocócica 10-valente, poliomielite, meningocócica C, rotavírus e primeira dose da tríplice viral. Entretanto, a recuperação não é homogênea entre os imunizantes, permanecendo níveis baixos em vacinas como DTP de primeiro reforço, hepatite A e segunda dose da tríplice viral, chamando atenção para o cuidado com as doses de reforço.

\vertabelas{Tabela~\ref{tab:cobertura-vacinal-epi-por-ano}; ver Apêndice~\ref{ap:eixo-3}}

\section{Famílias no CadÚnico com crianças até 72 meses, por renda e arranjo familiar}

\figura{grafico}{../../visualizacoes/a4/cadunico_familias_por_arranjo.pdf}{CADÚNICO: Famílias com crianças de 0 a 5 anos, por arranjo familiar}{Elaboração IPP com dados de \citeonline{mds_cadunico}.}{graf:cadunico-familias-por-arranjo}

Entre as famílias do CadÚnico com crianças de 0 a 5 anos, há forte predominância do arranjo com uma única adulta mulher, que reúne 145.643 famílias, cerca de 81,4\% do total. Em seguida aparecem as famílias com dois adultos, homem e mulher, com 18.489 registros, enquanto os demais arranjos possuem participação bem menor.

O resultado evidencia uma concentração expressiva de crianças cadastradas em famílias compostas por apenas uma mulher adulta. Entretanto, essa categoria não deve ser interpretada diretamente como família monoparental, pois o dado considera apenas a composição de adultos registrada no CadÚnico e não informa o parentesco entre os membros da família.

\figura{grafico}{../../visualizacoes/a4/cadunico_familias_arranjo_renda.pdf}{CADÚNICO: Renda per capita das famílias com crianças de 0 a 5 anos, por arranjo familiar}{Elaboração IPP com dados de \citeonline{mds_cadunico}.}{graf:cadunico-familias-arranjo-renda}

A renda per capita varia bastante conforme o arranjo familiar. Entre as famílias com uma única mulher adulta, 79,2\% estão na faixa de pobreza, enquanto apenas 5,2\% possuem renda acima de meio salário mínimo. Entre aquelas com dois adultos, homem e mulher, a proporção em situação de pobreza é menor, de 53,4\%, e 15,9\% estão acima de meio salário mínimo.

Também se destaca o grupo com um único homem adulto, no qual 72,8\% das famílias estão na faixa de pobreza. Já as famílias com três ou mais adultos apresentam distribuição menos concentrada, com 48,9\% em situação de pobreza e 37,8\% na faixa de baixa renda. De forma geral, os dados mostram maior concentração de pobreza nos arranjos com apenas um adulto.

\figura{mapa}{../../mapas/a4/mapa_percentual_cadunico_familias_uma_adulta_bairro_2026.pdf}{Famílias com crianças de 0 a 5 anos no CadÚnico: \% com uma só adulta, por bairro}{Elaboração IPP com dados de \citeonline{mds_cadunico}. Limites de bairros: \citeonline{ipp_limites_bairros}. Sistema de referência SIRGAS 2000.}{mapa:mapa-percentual-cadunico-familias-uma-adulta-bairro-2026}

Entre as famílias do CadÚnico com crianças de 0 a 5 anos, a presença de uma única adulta mulher é predominante em praticamente todo o município. A mediana entre os bairros é de aproximadamente 80\%, com valores que variam de 62,9\% em Vista Alegre a 88,1\% na Gamboa.

Os maiores percentuais aparecem também em bairros como Jacarezinho, Acari, Maré, Costa Barros e Manguinhos, todos acima de 85\%. O mapa mostra, portanto, que esse tipo de arranjo familiar não está restrito a poucos territórios, mas apresenta elevada participação em grande parte da cidade.

\vertabelas{Tabela~\ref{tab:cadunico-familias-por-arranjo-2026}, Tabela~\ref{tab:cadunico-familias-arranjo-renda-2026}, Tabela~\ref{tab:tabela-mapa-cadunico-recortes-bairro-2026}; ver Apêndice~\ref{ap:eixo-3}}

\section{Síntese do eixo}

\begin{sintese}{Família e Cuidados}Acreditamos que a demanda de pré-escola e escola será crescentemente concentrada na Zona Oeste da cidade. As taxas de matrícula abaixo dos objetivos apontam, entretanto, que o desafio não se resume a oferta, mas a capacidade de garantir a efetivação da matrícula por parte das famílias. Os arranjos familiares podem ser decorrentes de erros ou falsas declarações durante o cadastramento, mas exigem atenção quanto a possível demanda por ampliação das políticas de apoio a mães solteiras, entre as quais a oferta de creche e pré-escola em tempo integral.\end{sintese}

\eixo{4}{7}{shield-check}
